% ============================================================
%  Chapter 1 — What Is a Robot?
%  Robot World: Meet the Machines That Help Us
% ============================================================

\chapter{What Is a Robot?}

\chapteropening{%
  They build cars, explore volcanoes, vacuum floors, and even perform surgery.
  But what exactly \emph{is} a robot?
}
\chapterbadge{Mission Start: Enter Robot World!}

% --- Opening story ---
\section*{A New Adventure Begins}

Sam was visiting the local science museum. In the main hall, a
robotic arm was picking up coloured blocks and stacking them into a perfect
tower. Beside it, a small wheeled robot was drawing a picture on the floor
with a marker.

Sam stared, amazed:

\begin{center}
  \Large\bfseries\color{RoboGreen!80!black}
  ``How do they know what to do? Nobody is holding a controller!''
\end{center}

\vspace{8pt}

A familiar glow appeared on the screen of a nearby display.

\character{Dot}{Remember me? We have explored the Internet and the inside
  of a smartphone together. Now I want to show you something even more
  exciting \textemdash{} the world of robots!}

\sectionrule

% --- What makes a robot? ---
\section*{Three Things Every Robot Does}

Not every machine is a robot. A toaster heats bread, but it is not a robot.
So what makes something a robot? Every true robot does three things:

\begin{enumerate}
  \item \textbf{Sense} \textemdash{} it collects information about the world using \term{sensors}.
  \item \textbf{Think} \textemdash{} it processes that information using a \term{controller} (a kind of computer brain).
  \item \textbf{Act} \textemdash{} it does something in the real world using \term{actuators} (motors, wheels, grippers).
\end{enumerate}

\begin{bigidea}
  A robot is a machine that can \textbf{sense} the world, \textbf{think}
  about what to do, and then \textbf{act} on its own.
\end{bigidea}

\begin{picturethis}
  Imagine you are playing football. Your eyes \emph{sense} the ball coming.
  Your brain \emph{thinks}: ``I should kick it left.'' Your leg \emph{acts}
  and kicks. A robot works the same way \textemdash{} eyes, brain, body \textemdash{}
  just made of metal and code instead of flesh and bone.
\end{picturethis}

\sectionrule

% --- Robots vs machines ---
\section*{Robot or Just a Machine?}

Let us play a quick sorting game.

\begin{center}
\begin{tabularx}{0.85\textwidth}{|X|c|}
  \hline
  \textbf{Machine} & \textbf{Robot?} \\
  \hline
  A light switch & No \\
  \hline
  A self-driving car & Yes \\
  \hline
  A washing machine & No (it follows a fixed timer) \\
  \hline
  A robot vacuum that avoids furniture & Yes \\
  \hline
  A vending machine & No \\
  \hline
  A drone that follows you while filming & Yes \\
  \hline
\end{tabularx}
\end{center}

The key difference? Robots \emph{respond} to what is happening around them.
They make decisions.

\begin{funfact}
  The word ``robot'' comes from the Czech word \emph{robota}, meaning
  forced labour. It was first used in a 1920 play called \emph{R.U.R.}
  by Karel \v{C}apek.
\end{funfact}

\sectionrule

% --- Types of robots ---
\section*{Robots Come in All Shapes}

Robots do not all look like metal people. In fact, most robots look nothing
like humans at all!

\begin{itemize}
  \item \textbf{\term{Robot arms}} \textemdash{} bolted to factory floors, they weld, paint, and assemble.
  \item \textbf{\term{Wheeled robots}} \textemdash{} rovers, delivery bots, robot vacuums.
  \item \textbf{\term{Drones}} \textemdash{} flying robots used for filming, deliveries, and search-and-rescue.
  \item \textbf{\term{Humanoids}} \textemdash{} shaped like people, designed to work alongside us.
  \item \textbf{\term{Soft robots}} \textemdash{} made of flexible materials, inspired by octopuses and worms.
  \item \textbf{\term{Micro-bots}} \textemdash{} tiny robots that can travel inside the human body.
\end{itemize}

\character{Bolt}{I am a robot arm! I can lift heavy car doors all day long
  without getting tired. But I would be terrible at climbing stairs \textemdash{}
  that is a job for a different kind of robot.}

\begin{funfact}
  There are over 3.9 million industrial robots working in factories
  around the world --- that is roughly one robot for every 2,000 people
  on Earth. And the number is growing every year!
\end{funfact}

\sectionrule

% --- Try It ---
\begin{tryit}
  Look around your home. Can you find anything that senses, thinks, and acts?
  A robot vacuum? A smart speaker? Make a list and decide: is it a true
  robot, or just a clever machine?
\end{tryit}

\sectionrule

% --- Diagram ---
\begin{diagram}[The sense\,--\,think\,--\,act loop that every robot follows.]
  \begin{tikzpicture}[>=Stealth, node distance=3.5cm]
    \node[draw=RoboGreen!80!black,fill=RoboGreen!15,rounded corners=8pt,
          minimum width=2.5cm,minimum height=1.2cm,line width=1.2pt,
          font=\bfseries,align=center] (sense) {Sense};
    \node[draw=WarmOrange,fill=WarmOrange!15,rounded corners=8pt,
          minimum width=2.5cm,minimum height=1.2cm,line width=1.2pt,
          font=\bfseries,align=center,right=of sense] (think) {Think};
    \node[draw=SteelBlue,fill=SteelBlue!15,rounded corners=8pt,
          minimum width=2.5cm,minimum height=1.2cm,line width=1.2pt,
          font=\bfseries,align=center,right=of think] (act) {Act};
    \draw[->,line width=1.5pt,RoboGreen!70!black] (sense) -- (think);
    \draw[->,line width=1.5pt,WarmOrange!80!black] (think) -- (act);
    \draw[->,line width=1.5pt,SteelBlue!80!black,bend left=40] (act.south) to
      node[below,font=\small\itshape]{repeat} (sense.south);
  \end{tikzpicture}
\end{diagram}

\sectionrule


\begin{quickcheck}
  \begin{enumerate}
    \item What three things does every robot do?
    \item Name one job that is better for a robot than a human. Why?
    \item Can you find one robot (or robot-like machine) in your home?
  \end{enumerate}
\end{quickcheck}

\sectionrule
% --- New Words ---
\begin{newwords}
  \begin{description}[labelwidth=3.8cm, leftmargin=4.0cm]
    \item[Robot] A machine that can sense, think, and act on its own.
    \item[Sensor] A part that detects things (light, sound, distance, touch).
    \item[Controller] The computer brain that makes decisions.
    \item[Actuator] A part that moves or does work (motor, gripper, propeller).
    \item[Humanoid] A robot shaped like a human body.
    \item[Drone] An unmanned flying robot.
  \end{description}
\end{newwords}
